% Source: https://artofproblemsolving.com/community/c3h1885714_question_about_lims
% Stable identity: aops-topic-1885714-q01
(a) For $x\ge0$, $x\ne4$, multiply numerator and denominator by $\sqrt x+2$:
\[
\dfrac{\sqrt x-2}{x-4}=\dfrac{(\sqrt x-2)(\sqrt x+2)}{(x-4)(\sqrt x+2)}=\dfrac{x-4}{(x-4)(\sqrt x+2)}=\dfrac{1}{\sqrt x+2}\longrightarrow\dfrac{1}{4} .
\]
(b) For $x\ge0$, $x-4=\left(\sqrt x\right)^2-2^2=(\sqrt x-2)(\sqrt x+2)$, so again $\dfrac{\sqrt x-2}{x-4}=\dfrac{1}{\sqrt x+2}\to\dfrac{1}{4}$.

The limit is $\boxed{\dfrac{1}{4}}$. Both methods rest on the same identity $(\sqrt x-2)(\sqrt x+2)=x-4$: rationalising is factoring with square roots allowed. Direct substitution gives the form $\dfrac{0}{0}$, which only says that more work is needed; here the limit exists. So ``if factoring does not work, the limit does not exist'' is not a valid rule.
